\chapter{TEORI FREDHOLM}

\section{Alternatif Fredholm}

Pada tahun 1903, Erik Ivar Fredholm menerbitkan sebuah makalah rintisan tentang persamaan integral dalam jurnal
\emph{Acta Mathematica}. Di antara banyak hasil pentingnya terdapat teorema yang kini kita kenal sebagai
\emph{alternatif Fredholm}. Kita menyatakan suatu versi hasil ini dalam bahasa
yang tersedia bagi Fredholm pada masa itu.

\begin{prop}[Alternatif Fredholm I]\label{004011} Tetapkan $\lambda \in \C\setminus\{0\}$ dan misalkan $k$ suatu fungsi kontinu bernilai kompleks
pada persegi satuan $[0,1] \times [0,1]$. \textbf{Entah} persamaan-persamaan
 \index{alternatif Fredholm!versi I}%
tak homogen
 \begin{align}
   \lambda f(s) &- \int_0^1 k(s,t)f(t)\,dt = g(s)\label{004011i} \text{ dan}    \tag{1}\\
   \conj{\lambda} h(s) &- \int_0^1 \conj{k(t,s)}h(t)\,dt = j(s)\label{004011ii} \tag{2}
 \end{align}
memiliki solusi $f$ dan $h$ untuk setiap $g$ dan $j$ yang diberikan, masing-masing, dengan solusi-solusi tersebut
tunggal; dalam hal ini persamaan homogen yang bersesuaian
 \begin{align} \lambda f(s) &- \int_0^1 k(s,t)f(t)\,dt = 0\label{004011iii} \text{ dan} \tag{3}\\
        \conj{\lambda} h(s) &- \int_0^1 \conj{k(t,s)}h(t)\,dt = 0\label{004011iv} \tag{4}
 \end{align}
hanya memiliki solusi trivial; \textbf{---atau---}\\
persamaan homogen \eqref{004011iii} dan \eqref{004011iv} masing-masing memiliki
sejumlah hingga (dan tak nol) solusi bebas linear yang sama banyaknya, yaitu $f_1,\dots,f_n$ dan $h_1,\dots,h_n$,
secara berturut-turut; dalam hal ini persamaan tak homogen \eqref{004011i} dan
\eqref{004011ii} memiliki solusi jika dan hanya jika $g$ dan $j$ memenuhi
 \begin{align} \int_0^1 h_k(t)\conj{g(t)}\,dt &= 0\label{004011v} \text{ dan} \tag{5}\\
              \int_0^1 j(t) \conj{f_k(t)}\,dt &= 0\label{004011vi} \tag{6}
 \end{align}
untuk $k = 1, \dots, n$.
\end{prop}

Pada tahun 1906, David Hilbert telah menyadari bahwa pengintegralan itu sendiri sangat sedikit berkaitan dengan
kebenaran hasil ini. Ia menemukan bahwa yang penting adalah kekompakan
operator integral yang dihasilkan. (Pada awal abad ke-$20$
\emph{kekompakan} disebut \emph{kontinuitas lengkap}. Istilah \emph{ruang Hilbert}
telah digunakan sejak tahun 1911 sehubungan dengan ruang barisan~$l_2$. Baru pada
tahun 1929 John von Neumann memperkenalkan gagasan---dan aksioma yang
mendefinisikan---\emph{ruang Hilbert abstrak}.) Jadi, berikut suatu versi yang agak lebih mutakhir dari
\emph{alternatif Fredholm}.

\begin{prop}[Alternatif Fredholm II]\label{004031} Jika $K$ suatu operator kompak pada ruang
 \index{alternatif Fredholm!versi II}%
Hilbert, $\lambda \in \C\setminus\{0\}$, dan $T = \lambda I - K$, maka \textbf{entah} persamaan-persamaan
tak homogen
 \begin{align}
            Tf &= g\label{004031i} \text{ dan} \tag{1'}\\
            T^* h &= j\label{004031ii} \tag{2'}
 \end{align}
memiliki solusi $f$ dan $h$ untuk setiap $g$ dan $j$ yang diberikan, masing-masing, dengan solusi-solusi tersebut
tunggal; dalam hal ini persamaan homogen yang bersesuaian
 \begin{align}
            Tf &= 0\label{004031iii} \text{ dan} \tag{3'}\\
          T^*h &= 0\label{004031iv} \tag{4'}
 \end{align}
hanya memiliki solusi trivial; \textbf{---atau---}\\
persamaan homogen \eqref{004031iii} dan \eqref{004031iv} masing-masing memiliki
sejumlah hingga (dan tak nol) solusi bebas linear yang sama banyaknya, yaitu $f_1,\dots,f_n$ dan $h_1,\dots,h_n$,
secara berturut-turut; dalam hal ini persamaan tak homogen \eqref{004031i} dan
\eqref{004031ii} memiliki solusi jika dan hanya jika $g$ dan $j$ memenuhi
 \begin{align} h_k &\perp g\label{004031v} \text{ dan} \tag{5'}\\
                 j &\perp f_k\label{004031vi} \tag{6'}
 \end{align}
untuk $k = 1, \dots, n$.
\end{prop}

Perhatikan bahwa dengan menggunakan beberapa fakta elementer mengenai kernel dan jangkauan
operator serta keortogonalan dalam ruang Hilbert, kita dapat meringkas pernyataan
dari~\ref{004031} secara cukup banyak.

\begin{prop}[Alternatif Fredholm IIIa]\label{004033} Jika $T = \lambda I - K$ dengan
 \index{alternatif Fredholm!versi IIIa}%
$K$ suatu operator kompak pada ruang Hilbert dan $\lambda \in \C\setminus\{0\}$, maka
 \begin{enumerate}
  \item $T$ injektif jika dan hanya jika bersifat surjektif,
  \item $\ran T^* = (\ker T)^\perp$, dan
  \item $\dim\ker T = \dim \ker T^*$.
 \end{enumerate}
Selain itu, syarat (1) dan (2) berlaku untuk $T^*$ maupun~$T$.
\end{prop}


















\section{Alternatif Fredholm -- lanjutan}
\begin{defn}\label{004034} Suatu operator $T$ pada ruang Banach disebut
 \index{operator Riesz--Schauder}%
 \index{operator!Riesz--Schauder}%
\df{operator Riesz--Schauder} jika dapat ditulis dalam bentuk $T = S + K$ dengan $S$
invertibel dan $K$ kompak.
\end{defn}

Materi dalam bagian~\ref{sec_onbases} dan definisi sebelumnya memungkinkan kita
memperumum versi \emph{alternatif Fredholm} yang diberikan dalam~\ref{004033} ke ruang
Banach.

\begin{prop}[Alternatif Fredholm IIIb]\label{0040343} Jika $T$ suatu operator Riesz--Schauder
 \index{alternatif Fredholm!versi IIIb}%
pada ruang Banach, maka
 \begin{enumerate}
  \item $T$ injektif jika dan hanya jika bersifat surjektif,
  \item $\ran T^* = (\ker T)^\perp$, dan
  \item $\dim\ker T = \dim \ker T^* < \infty$.
 \end{enumerate}
Selain itu, syarat (1) dan (2) berlaku untuk $T^*$ maupun~$T$.
\end{prop}

\begin{prop} Jika $M$ subruang tertutup dari ruang Banach $B$, maka $M^\perp \cong (B/M)^*$.
\end{prop}

\begin{proof} Lihat \cite{Conway:1990}, halaman 89. \ns \end{proof}

\begin{defn} Jika $T\colon V \sto W$ suatu pemetaan linear antara ruang vektor, maka
 \index{kokernel}%
\df{kokernel} pemetaan tersebut didefinisikan oleh
  \[ \coker T = W/ \ran T\,. \]
Ingat bahwa
 \index{kodimensi}%
\df{kodimensi} dari suatu subruang $U$ dalam ruang vektor $V$ adalah dimensi dari~$V/U$. Jadi,
apabila $T$ suatu pemetaan linear, $\dim\coker T = \codim\ran T$.
\end{defn}

Dalam kategori $\cat{HIL}$ dari ruang Hilbert dan pemetaan linear kontinu, jangkauan suatu
morfisme tidak harus menjadi objek kategori tersebut. Secara khusus, jangkauan suatu operator
tidak harus tertutup.

\begin{exam}\label{004066} Operator ruang Hilbert
   \[ T\colon l_2 \sto l_2\colon (x_1,x_2,x_3,\dots) \mapsto
                                   (x_1, \tfrac12 x_2, \tfrac13 x_3, \dots) \]
bersifat injektif, swaadjoin, kompak, dan kontraktif,
 \index{operator!dengan jangkauan tak tertutup}%
tetapi jangkauannya, walaupun padat dalam $l_2$, tidak mencakup seluruh~$l_2$.
\end{exam}

Walaupun tidak langsung berkaitan dengan tujuan kita saat ini, inilah tempat yang tepat untuk mencatat
fakta bahwa jumlah dua subruang tertutup dari suatu ruang Hilbert tidak harus tertutup.

\begin{exam}\label{004071} Misalkan $T$ operator pada ruang Hilbert $l_2$ yang didefinisikan dalam
 \index{subruang!jumlah subruang tertutup tidak harus tertutup}%
contoh~\ref{004066}, $M = l_2 \oplus \{\vc 0\}$, dan $N$ graf dari~$T$. Maka $M$
dan $N$ keduanya merupakan subruang tertutup dari ruang Hilbert $l_2 \oplus l_2$, tetapi $M + N$
tidak tertutup.
\end{exam}

\begin{proof} Verifikasi contoh~\ref{004071} diperoleh dengan mudah dari hasil berikut
dan contoh~\ref{004066}. \ns
\end{proof}

\begin{prop} Misalkan $H$ suatu ruang Hilbert, $T \in \ofml B(H)$, $M = H \oplus \{\vc 0\}$, dan $N$
merupakan graf dari~$T$. Maka
 \begin{enumerate}
   \item[(a)] Himpunan $N$ merupakan subruang dari $H \oplus H$.
   \item[(b)] Operator $T$ injektif jika dan hanya jika $M \cap N = \{\,(\vc 0,\vc 0)\,\}$.
   \item[(c)] Jangkauan $T$ padat dalam $H$ jika dan hanya jika $M + N$ padat dalam $H \oplus H$.
   \item[(d)] Operator $T$ surjektif jika dan hanya jika $M + N = H \oplus H$.
 \end{enumerate}
\end{prop}

Kabar baik bagi teori operator Fredholm adalah bahwa operator dengan kokernel
berdimensi hingga secara otomatis mempunyai jangkauan tertutup.

\begin{prop}\label{00411} Jika suatu pemetaan linear terbatas $A\colon H \sto K$ antara ruang Hilbert
memiliki kokernel berdimensi hingga, maka $\ran A$ tertutup dalam~$K$.
\end{prop}

Dalam \emph{Alternatif Fredholm IIIb}~\ref{0040343}, kita mengamati bahwa syarat (1)
berlebih dan juga bahwa (2) berlaku untuk \emph{sebarang} operator ruang Banach dengan jangkauan tertutup.
Hal ini memungkinkan kita menyatakan kembali~\ref{0040343} secara lebih hemat.

\begin{prop}[Alternatif Fredholm IV]\label{00413}
 \index{alternatif Fredholm!versi IV}%
Jika $T$ suatu operator Riesz--Schauder pada ruang Banach, maka
 \begin{enumerate}
    \item $T$ mempunyai jangkauan tertutup dan
    \item  $\dim\ker T = \dim \ker T^* < \infty$.
 \end{enumerate}
\end{prop}














\section{Operator Fredholm}
\begin{defn} Misalkan $H$ adalah ruang Hilbert dan $\ofml K(H)$ merupakan ideal operator kompak di dalam
aljabar-$C^*$~$\ofml B(H)$. Maka aljabar hasil bagi (lihat proposisi~\ref{001902})
$\ofml Q(H) := \ofml B(H)/\ofml K(H)$ adalah
 \index{aljabar Calkin}%
 \index{aljabar!Calkin}%
 \index{Q@$\ofml Q(H)$ (aljabar Calkin)}%
\df{aljabar Calkin}. Seperti biasa, pemetaan hasil bagi dari $\ofml B(H)$ ke $\ofml
B(H)/\ofml K(H)$
 \index{pi@$\pi$ (pemetaan hasil bagi)}%
 \index{pi@$\pi(T)$ (elemen aljabar Calkin yang memuat $T$)}%
dilambangkan dengan $\pi$, sehingga jika $T \in \ofml B(H)$ maka $\pi(T) = [T] = T + \ofml K(H)$ adalah
elemen yang bersesuaian dengannya dalam aljabar Calkin. Suatu elemen $T \in \ofml B(H)$ adalah
 \index{Fredholm!operator}%
 \index{operator!Fredholm}%
\df{operator Fredholm} jika $\pi(T)$ invertibel di dalam~$\ofml Q(H)$. Keluarga
 \index{F@$\ofml F(H)$ (operator Fredholm pada $H$)}%
semua operator Fredholm pada $H$ dilambangkan dengan~$\ofml F(H)$.
\end{defn}

\begin{prop} Jika $H$ ruang Hilbert, maka $\ofml F(H)$ merupakan subhimpunan terbuka swaadjoin
dari~$\ofml B(H)$ dan tertutup terhadap
 \index{perturbasi}%
perturbasi kompak (yakni, jika $T$ Fredholm dan $K$ kompak, maka $T + K$ juga Fredholm).
\end{prop}

\begin{thm}[Teorema Atkinson]\label{004352}
 \index{teorema Atkinson}%
Suatu operator ruang Hilbert bersifat Fredholm jika dan hanya jika kernel dan kokernelnya
berdimensi hingga.
\end{thm}

\begin{proof} Lihat \cite{Arveson:2002}, teorema 3.3.2; \cite{Blackadar:2006}, teorema I.8.3.6;
\cite{Douglas:1972}, teorema 5.17; \cite{HigsonR:2000}, teorema 2.1.4;
atau~\cite{Wegge-Olsen:1993}, teorema 14.1.1. \ns
\end{proof}

Menariknya, dimensi kernel dan kokernel suatu operator Fredholm bukanlah besaran yang
sangat penting. Namun, selisih keduanya ternyata sangat penting.

\begin{defn} Jika $T$ adalah operator Fredholm pada suatu ruang Hilbert, maka
 \index{indeks!operator Fredholm}%
 \index{Fredholm!indeks (\seeonly{indeks})}%
\df{indeks Fredholm} operator tersebut (atau cukup \df{indeks}) didefinisikan oleh
   \[ \ind T := \dim\ker T - \dim\coker T\,. \]
\end{defn}

\begin{exam}\label{0044125} Setiap operator ruang Hilbert yang invertibel bersifat Fredholm
 \index{indeks!operator invertibel}%
 \index{operator!invertibel!indeks}%
 \index{invertibel!operator!indeks}%
dengan indeks nol.
\end{exam}

\begin{exam}\label{0044128} Indeks Fredholm dari
 \index{indeks!operator geser unilateral}%
 \index{operator geser unilateral!indeks}%
 \index{operator!geser unilateral!indeks}%
operator geser unilateral adalah~$-1$.
\end{exam}

\begin{exam} Untuk pemetaan antar-ruang, kita memakai konvensi standar: pemetaan linear disebut Fredholm jika
jangkauannya tertutup serta kernel dan kokernelnya berdimensi hingga. Dengan konvensi ini, setiap pemetaan linear $T\colon V \sto W$ antara ruang-ruang vektor berdimensi hingga bersifat Fredholm dan
$\ind T = \dim V - \dim W$. \end{exam}

\begin{exam}\label{004413} Jika $T$ adalah operator Fredholm pada suatu ruang Hilbert,
 \index{indeks!adjoin operator}%
 \index{operator!adjoin!indeks}%
 \index{adjoin!indeks}%
maka $\ind T^* = -\ind T$.
\end{exam}

\begin{exam} Indeks setiap operator Fredholm
 \index{indeks!operator Fredholm normal}%
 \index{operator!Fredholm normal!indeks}%
 \index{normal!operator Fredholm!indeks}%
normal adalah~$0$.
\end{exam}

\begin{lem}\label{004432} Misalkan $S\colon U \sto V$ dan $T\colon V \sto W$ adalah transformasi linear
antara ruang-ruang vektor. Maka (terdapat pemetaan linear sedemikian sehingga) barisan
berikut eksak.
 \[ \vc 0 \to \ker S \to \ker TS \to \ker T \to \coker S
                            \to \coker TS \to \coker T \to \vc 0\,. \]
\end{lem}

\begin{lem}\label{004433} Jika $V_0, V_1, \dots, V_n$ adalah ruang-ruang vektor berdimensi hingga
dan barisan
  \[ \vc 0 \to V_0 \to V_1 \to \dots \to V_n \to \vc 0 \]
eksak, maka
  \[ \sum_{k=0}^n (-1)^k \dim V_k = 0\,. \]
\end{lem}

\begin{prop} Misalkan $H$ ruang Hilbert berdimensi tak hingga. Maka himpunan $\ofml F(H)$ yang terdiri atas operator Fredholm
pada $H$ merupakan semigrup terhadap komposisi dan fungsi indeks $\ind$ adalah epimorfisme
dari $\ofml F(H)$ ke semigrup aditif bilangan bulat~$\Z$.
\end{prop}

\begin{proof} \emph{Petunjuk.} Gunakan \ref{004432}, \ref{004433}, \ref{0044125}, \ref{0044128},
dan~\ref{004413}. \ns
\end{proof}





















\section{Alternatif Fredholm -- Penutup}
Hasil berikut menyiratkan bahwa setiap operator Fredholm berindeks nol dapat ditulis sebagai
jumlah suatu operator invertibel dan operator kompak.

\begin{prop}\label{004712} Jika $T$ adalah operator Fredholm berindeks nol pada suatu ruang Hilbert,
maka terdapat isometri parsial berperingkat hingga $V$ sedemikian sehingga $T - V$ bersifat invertibel.
\end{prop}

\begin{proof} Lihat \cite{Pedersen:1995}, lemma 3.3.14 atau \cite{Wegge-Olsen:1993}, proposisi
14.1.3.  \ns
\end{proof}

\begin{lem} Jika $F$ adalah operator berperingkat hingga pada suatu ruang Hilbert, maka $I + F$ bersifat Fredholm dengan
indeks nol.
\end{lem}

\begin{proof} Lihat \cite{Pedersen:1995}, lemma 3.3.13 atau \cite{Wegge-Olsen:1993}, lemma 14.1.4. \ns
\end{proof}

\begin{notn} Misalkan $H$ adalah ruang Hilbert. Untuk setiap bilangan bulat $n$, kita melambangkan
 \index{F@$\ofml F_n(H)$, $\ofml F_n$ (operator Fredholm berindeks~$n$)}%
keluarga semua operator Fredholm berindeks $n$ pada $H$ dengan $\ofml F_n(H)$ atau cukup $\ofml
F_n$.
\end{notn}

\begin{prop}\label{004724} Untuk setiap ruang Hilbert berlaku $\ofml F_0 + \ofml K = \ofml F_0$.
\end{prop}

\begin{proof} Lihat \cite{Douglas:1972}, lemma 5.20 atau \cite{Wegge-Olsen:1993}, 14.1.5. \ns
\end{proof}

\begin{cor}[Alternatif Fredholm V]\label{004725}
 \index{alternatif Fredholm!versi V}%
Setiap operator Riesz--Schauder pada suatu ruang Hilbert bersifat Fredholm berindeks nol.
\end{cor}

\begin{proof} Ini segera mengikuti proposisi sebelumnya~\ref{004724} dan
contoh~\ref{0044125}.
\end{proof}

Kita sebenarnya telah membuktikan hasil yang lebih kuat: versi terakhir (yang cukup umum) dari
\emph{alternatif Fredholm}.

\begin{cor}[Alternatif Fredholm VI]\label{0047253} Suatu operator ruang Hilbert merupakan operator
 \index{alternatif Fredholm!versi VI}%
Riesz--Schauder jika dan hanya jika operator itu merupakan operator Fredholm berindeks nol.
\end{cor}

\begin{proof} Proposisi~\ref{004712} dan korolari~\ref{004725}.
\end{proof}

\begin{prop} Jika $T$ adalah operator Fredholm berindeks $n \in \Z$ pada ruang Hilbert $H$
dan $K$ operator kompak pada~$H$, maka $T + K$ juga merupakan operator Fredholm berindeks~$n$; yakni,
  \[ \ofml F_n + \ofml K = \ofml F_n\,. \]
\end{prop}

\begin{proof} Lihat \cite{HigsonR:2000}, proposisi 2.1.6; \cite{Pedersen:1995}, teorema
3.3.17; atau \cite{Wegge-Olsen:1993}, proposisi 14.1.6. \ns
\end{proof}

\begin{prop} Misalkan $H$ adalah ruang Hilbert. Maka $\ofml F_n(H)$ merupakan subhimpunan terbuka dari $\ofml B(H)$
untuk setiap bilangan bulat~$n$.
\end{prop}

\begin{proof} Lihat \cite{Wegge-Olsen:1993}, proposisi 14.1.8. \ns \end{proof}

\begin{defn}\label{004735} Suatu
 \index{lintasan}%
\df{lintasan} dalam ruang topologis $X$ adalah pemetaan kontinu dari interval $[0,1]$
ke~$X$. Dua titik $p$ dan $q$ dalam $X$ dikatakan
 \index{lintasan!titik-titik yang terhubung oleh}%
 \index{terhubung!oleh lintasan}%
\df{terhubung oleh lintasan} (atau
 \index{homotop!titik-titik dalam ruang topologis}%
\df{homotop dalam~$X$}) jika terdapat lintasan $f\colon [0,1] \sto X$ dalam $X$ sedemikian sehingga
$f(0) = p$ dan $f(1) = q$. Dalam hal ini kita
 \index{<binrelequivh@$p \sim_h q$ (homotopi titik-titik)}%
menulis $p \sim_h q$ dalam~$X$ (atau cukup $p \sim_h q$ ketika ruang $X$ jelas dari konteks).
\end{defn}


\begin{exam} Misalkan $a$ adalah elemen invertibel dalam aljabar-$C^*$ beridentitas $A$ dan $b$ adalah elemen $A$
sedemikian sehingga $\norm{a - b} < \norm{a^{-1}}^{-1}$. Maka $a \sim_h b$ dalam~$\inv(A)$.
\end{exam}

\begin{proof}[\emph{Petunjuk bukti}] Terapkan korolari~\ref{cor2_Neumann_series} pada titik-titik dalam ruas garis tertutup~$[a,b]$. \ns
\end{proof}

\begin{prop}\label{K002011} Relasi $\sim_h$ untuk ekuivalensi homotopi yang didefinisikan di atas merupakan suatu relasi
 \index{ekuivalensi!homotopi!titik-titik}%
ekuivalensi pada himpunan titik-titik suatu ruang topologis.
\end{prop}

\begin{defn} Jika $X$ adalah ruang topologis dan $\sim_h$ adalah relasi ekuivalensi homotopi,
maka kelas-kelas ekuivalensi yang dihasilkan adalah
 \index{lintasan!komponen}%
 \index{komponen!lintasan}%
\df{komponen lintasan} dari~$X$.
\end{defn}

Proposisi berikut mengidentifikasi komponen-komponen lintasan himpunan operator Fredholm sebagai
himpunan $\ofml F_n$ yang terdiri atas operator berindeks~$n$.

\begin{prop} Operator $S$ dan $T$ di dalam ruang operator Fredholm $\ofml F(H)$ pada
ruang Hilbert $H$ bersifat homotop dalam $\ofml F(H)$ jika dan hanya jika keduanya memiliki indeks yang sama.
\end{prop}

\begin{proof} Lihat \cite{Wegge-Olsen:1993}, korolari 14.1.9. \ns \end{proof}







\endinput
